\section{Stage IX: Explicit SIMD in OxCaml}
\label{sec:simd}

The last remaining difference was vectorisation, so I added an OxCaml kernel
that uses SIMD explicitly (build switch \code{MG_SIMD=1}). Splitting one row's
dot product across lanes would reassociate the sum and break the bit-exact
contract, so the kernel puts \emph{rows} in lanes, each adding the same
products in the same order as the scalar kernel, as GCC's vectoriser does.

\begin{lstlisting}[language=OCamlX,caption={The SSE2 kernel's inner loop. \texttt{F.add} and \texttt{F.mul} are \texttt{[@@builtin]} externals on the unboxed 128-bit type \texttt{float64x2\#}.},label=lst:simd]
let rec go i (v01 : float64x2#) (v23 : float64x2#) =
  if i >= cols then #(v01, v23)
  else
    let xi = unbox x.(i) in
    let xv = F.pair xi xi in                               (* broadcast x.(i) *)
    let r01 = F.pair (unbox d.(b0 + i)) (unbox d.(b1 + i)) in  (* rows 0, 1 *)
    let r23 = F.pair (unbox d.(b2 + i)) (unbox d.(b3 + i)) in  (* rows 2, 3 *)
    go (i + 1) (F.add v01 (F.mul r01 xv)) (F.add v23 (F.mul r23 xv))
\end{lstlisting}

The intrinsics are expanded inline, so no C call happens at run time. All
SIMD builds produce the reference output.

\subsection{Results: SIMD Helps Most Where the Passes Cannot Act}
\label{fig:interaction}
GCC's vectoriser is worth @P{cpp_novec,cpp_unroll}@ to C++. On the
pull-request compiler with the passes off, the SSE2 kernel is worth
@P{pr_off,pr_simd_off}@. With the passes on it adds much less:
@P{pr_all,pr_simd}@ (interquartile range @I{pr_all,pr_simd}@) with a poll on
every back edge, @P{pr_nopoll,pr_simd_nopoll}@ with no polls, and
@P{pr_counted,pr_simd_counted}@ with counted loops, leaving OxCaml with the
SSE2 kernel @G{pr_simd_counted}@ behind scalar C++, against @G{pr_counted}@
for the scalar kernel. SIMD and the pointer passes remove
overlapping work.

\subsection{Why SIMD Adds Little to the Passes}
\label{sec:simd-why}
Interior pointers and the shared row index survive in the SIMD kernel
(Listing~\ref{lst:simd-asm}); without the passes it needs four more
\code{lea}s per iteration (20 instructions instead of 16).

\begin{lstlisting}[language=asmX,caption={The SSE2 kernel's inner loop without safepoints: 16 instructions per iteration against the scalar loop's 13, the extra work being packing.},label=lst:simd-asm]
vmovsd   -0x4(%rax,%rbp,4),%xmm2     # x[i]
vmovapd  %xmm2,%xmm3
vmovlhps %xmm2,%xmm3,%xmm2           # broadcast x[i] into both lanes
vmovsd   -0x8(%r10,%rbp,4),%xmm3     # row 0: interior pointer + shared index
vmovsd   -0x8(%r13,%rbp,4),%xmm4     # row 1
vmovlhps %xmm4,%xmm3,%xmm3           # pack (row 0, row 1)
vmovsd   -0x8(%r12,%rbp,4),%xmm4     # row 2
vmovsd   -0x8(%r9,%rbp,4),%xmm5      # row 3
vmovlhps %xmm5,%xmm4,%xmm4           # pack (row 2, row 3)
vmulpd   %xmm2,%xmm4,%xmm4
vaddpd   %xmm4,%xmm0,%xmm0           # accumulators for rows 2, 3
vmulpd   %xmm2,%xmm3,%xmm2
vaddpd   %xmm2,%xmm1,%xmm1           # accumulators for rows 0, 1
add      $0x2,%rbp
cmp      %rbx,%rbp
jl       .loop
\end{lstlisting}

One row's elements are not adjacent to the next row's, so the kernel issues
the same five scalar loads and spends three shuffles and a move packing them,
in exchange for halving the floating-point instructions. It also has only two
accumulators where the scalar kernel has four independent chains.

If each loop is limited by its chains of dependent additions, doubling the
chains should speed up both kernels; if by the floating-point ports, extra
chains change nothing for the scalar kernel, while the SSE2 kernel, with half
as many operations, may be chain-limited. An eight-row version of both
kernels tests this (a separate 21-round run, identical output). Extra rows
barely help the scalar kernels ($+3.8\%$ for C++, nothing for OCaml), as a
port-limited loop predicts. The SSE2 kernel, now with four vector accumulators, gains $+5.1\%$
over its four-row version and beats the eight-row scalar kernel by $+4.7\%$
(21/21 rounds, with counted loops), where at four rows the two were equal in
that run. The
scalar loop is therefore port-limited, as the model of
Section~\ref{sec:residual} also finds, and the four-row SSE2 kernel
chain-limited.

In the counted build's hot copy of the SIMD loop, the fallback copy's live
values push three of the four interior pointers onto the stack (19
instructions): SIMD and interior pointers compete for registers. SIMD is the right tool where
the passes cannot act; with counted loops it adds a few percent, and the
scalar kernel, simpler and free of OxCaml extensions, gets most of the way.
